\chapter{Limits of the mapping}
\label{cap:limiti}

The ERA5 quantities of \cref{tab:limiti} have no counterpart in the archive. The limit lies
in the wrfout files, which do not contain a field computed by the model, in the configuration of the
simulation (\cref{sec:configurazione}) or in the publication rule (\cref{sec:pubblicazione}).

\begin{table}[htbp]
  \centering
  \small
  \caption{ERA5 quantities absent from the archive and origin of the limit.}
  \label{tab:limiti}
  \begin{tabularx}{\textwidth}{@{}XX@{}}
    \toprule
    Parameter & Origin of the limit \\
    \midrule
    boundary layer height (\campo{blh}) & wrfout files: \campo{PBLH} computed and not written \\
    sensible and latent heat fluxes & wrfout files: \campo{HFX} and \campo{LH} computed and not written \\
    turbulent kinetic energy & configuration: no scheme with turbulent kinetic energy \\
    maximum and minimum \SI{2}{m} temperature (\campo{mx2t}, \campo{mn2t}) & configuration: \campo{output_diagnostics=0} \\
    waves, sea-state-dependent drag coefficient & configuration: no wave model \\
    currents, sea surface height & configuration: ocean not simulated \\
    sea-ice thickness and albedo & configuration: sea ice only as input data \\
    convective precipitation (\campo{cp}, \campo{csf}) & configuration: resolved convection, identically zero field \\
    gust components (\campo{10efg}, \campo{10nfg}) & publication rule \\
    leaf area index per class (\campo{lai_lv}, \campo{lai_hv}) & publication rule \\
    angle and anisotropy of sub-grid orography (\campo{anor}, \campo{isor}) & publication rule \\
    snowmelt (\campo{smlt}) & publication rule \\
    the 8 ERA5 pressure levels above \SI{50}{\hPa} (\SIlist{1;2;3;5;7;10;20;30}{\hPa}) & configuration: model top at \SI{50}{\hPa} \\
    \bottomrule
  \end{tabularx}
\end{table}

\section{Fields computed and not written}
\label{sec:non-scritti}

The YSU scheme computes a boundary layer height (\campo{PBLH}), and with \campo{isfflx=1} the sensible
and latent heat fluxes (\campo{HFX}, \campo{LH}) are computed at every time step. None of the three
was written to the wrfout files, and no more complete output of the simulation exists. The same
holds for the model skin temperature \campo{TSK}, which the archive reconstructs from the radiation
(\cref{sec:pelle}).

For the boundary layer height, an a posteriori diagnosis, for example with the bulk Richardson number
on the profiles, would be ours and not the model's, and does not satisfy the publication rule
(\cref{sec:pubblicazione}). Nor can the fluxes be reconstructed from the surface energy balance,
because the balance keeps three unknowns: the two fluxes and the ground heat flux, which was not
written either.

\section{Fields excluded by the configuration}
\label{sec:configurazione-limiti}

The YSU scheme is non-local and has no prognostic turbulent kinetic energy, and with \campo{KM_OPT=4} the
diffusion does not produce it either: the field does not exist at any height.

With \campo{output_diagnostics=0} WRF did not write the hourly maxima and minima of the \SI{2}{m}
temperature, and an hourly sample of \campo{2t} is not the extreme within the hour. With
\campo{nwp_diagnostics=1} the simulation writes hourly maxima, but none for temperature.

No wave model is coupled to the atmosphere, so the significant wave height, the wave period and
direction, and the sea-state-dependent drag coefficient are missing. The ocean is not simulated
(\campo{SF_OCEAN_PHYSICS=0}), so there are no currents and no sea surface height. Sea ice comes only
from the input data, as a 0/1 indicator per cell, with a constant albedo of \num{0.65}; neither
the ice thickness nor an evolving ice albedo exists.

With resolved convection (\campo{CU_PHYSICS=0}) the convective precipitation \campo{RAINC} is
identically zero, and the ERA5 convective parameters \campo{cp} and \campo{csf} are not published
(\cref{sec:non-pubblicati}).

\section{Fields excluded by the publication rule}
\label{sec:esclusi}

\paragraph{Gust components.} For \campo{10efg} and \campo{10nfg} only the magnitude of the hourly
maximum of the resolved wind exists (\cref{cap:raffica}), and the two components are excluded by
decision.

\paragraph{Leaf area index.} The model carries a single leaf area index per cell, a table value for
the dominant category (\campo{rdlai2d=.false.}), whereas ERA5 requires one for low vegetation and one
for high vegetation. Any split of the single value between the two would be an assumption of ours.
The model field \campo{LAI} is in the wrfout files and is not published.

\paragraph{Angle and anisotropy of sub-grid orography.} The static file carries the asymmetry
(\campo{OA1-4}) and effective length (\campo{OL1-4}) parameters of the Kim--Arakawa scheme. ECMWF's
\campo{anor} and \campo{isor} derive from the structure of the covariance of the terrain gradient,
and the two quantities are different.

\paragraph{Snowmelt.} The field \campo{ACSNOM} exists but is not a monotonic accumulator. On 15 of the
16 days examined some points decrease between 00 and 12 UTC, by up to \SI{8.2}{mm}, and on 1 July
2024 the decrease relative to 00 UTC reaches \SI{7.26}{mm}. Built like the other accumulated
fields, as the difference between two consecutive hours of the same run (\cref{sec:tempo}),
\campo{smlt} would be negative at those points.

\section{Fields present and not published}
\label{sec:non-pubblicati}

Sixteen fields of the wrfout files are identically zero at every validity time of the verification sample:
\campo{ACHFX}, \campo{ACLHF}, \campo{ACGRDFLX}, \campo{NOAHRES}, \campo{SSTSK}, \campo{SST_INPUT},
\campo{SWNORM}, \campo{ACSNOW}, \campo{HAILNC}, \campo{RAINC}, \campo{RAINSH}, \campo{PREC_ACC_C} and
the four fields of downward thermal radiation at the top of the atmosphere, \campo{LWDNT},
\campo{LWDNTC}, \campo{ACLWDNT} and \campo{ACLWDNTC}. \campo{ACHFX}, \campo{ACLHF}, \campo{ACGRDFLX}
and \campo{NOAHRES} are zero because they are accumulators of the Noah surface scheme, and this
simulation uses RUC. An identically zero field is not published, whatever the reason.

For other non-zero fields no ECMWF parameter describes the same quantity: the frozen fraction of
precipitation (\campo{SR}), the density of falling snow (\campo{RHOSNF}), the annual maximum and
minimum of the green vegetation fraction (\campo{SHDMAX}, \campo{SHDMIN}) and the deep soil
temperature used as lower boundary condition (\campo{TMN}). The accumulated runoff \campo{ACRUNOFF}
coincides with the surface runoff \campo{SFROFF}, and is already in \campo{sro}.

Five non-zero fields stay out even though ECMWF defines a parameter for the same quantity, because
ERA5 does not publish it and the archive follows the ERA5 schema (\cref{tab:non-pubblicati}).

\begin{table}[htbp]
  \centering
  \small
  \caption{Non-zero fields of the wrfout files with an ECMWF parameter that ERA5 does not publish.}
  \label{tab:non-pubblicati}
  \begin{tabularx}{\textwidth}{@{}lX@{}}
    \toprule
    WRF field & ECMWF parameter \\
    \midrule
    \campo{SNOALB}     & 260161 \campo{mxsalb}, maximum snow albedo \\
    \campo{QGRAUP}     & 260028 \campo{grle} and 260001 \campo{tcolg}, graupel and its column (\cref{sec:colonne}) \\
    \campo{REFL_10CM}  & 260134 \campo{bref}, radar reflectivity \\
    \campo{LPI}        & family 228050 \campo{litoti}; \campo{LPI}, however, is a lightning potential index, not a flash density \\
    \campo{SNOWFALLAC} & 3066 \campo{sde}, because it is a geometric snow depth and not a water equivalent \\
    \bottomrule
  \end{tabularx}
\end{table}
